% Page 021 — page imprimée 17
% Suite du mémoire de Tandon : Une épidémie à Meyrueis en 1768.


{\fontsize{10.2}{12.0}\selectfont

Dans des temps plus heureux on les a vus se nourrir d’assez bon pain de froment, de bons légumes, d’un peu de viande fraîche, de cochon salé, de harengs et de morue, de châtaignes et d’assez bon fromage et se procurer une quantité suffisante de linge pour se tenir propres, et se coucher sur la paille qu’ils avaient soin de renouveler assez souvent….

Cette maladie, quoique très dangereuse par elle-même n’eut point eu de suites si fâcheuses si, par une fatalité dont on a peu d’exemples, la plupart des malades n’eussent été presque entièrement privés de secours :

Mr Auzilion, chirurgien de Meyrueis, qu’on a eu le malheur de perdre vers le milieu de cette épidémie a été pendant longtemps le seul qui les ait visités et leur ait administré des remèdes. Mais malgré sa bonne volonté et une très grande activité il lui était impossible de les visiter tous et aussi souvent qu’il eut fallu, vu leur nombre, leur éloignement et leur dispersion….

Depuis sa mort jusqu’à mon arrivée, aucun d’eux n’avait été vu que par M. Fignel curé de Lanuéjols..

Aussi en est-il mort un très grand nombre qui n’avaient pris pour tout remède que du suc d’oignon blanc, mêlé avec de la mauvaise huile ou de la suie détrempée dans du vinaigre ; remède réputé souverain parmi eux pour les maladies de la nature de celle-ci. qu’ils qualifiaient de peste blanche….

\medskip

\textit{Pour arrêter les progrès de cette maladie :}

\medskip

\textbf{1/} On exhortera tous les habitants de la campagne à se tenir aussi proprement qu’il serait possible en leur faisant sentir toute la nécessité d’emporter à des distances assez éloignées le fumier qui entoure leurs chaumières, de nettoyer le dedans de renouveler souvent leurs grabats, de changer de linge et de ne s’exposer au grand air que lorsque les brouillards seront entièrement dissipés.

\medskip

\textbf{2/} Il sera journellement distribué aux nécessiteux, en danger de tomber malades, à raison de leur détresse, un peu de viande fraîche, du cochon salé, du pain second, quelques légumes et même s’il était possible, un peu de vin et une certaine quantité de linge de rechange

\medskip

\textbf{3/} Les médecins feront tous les jours la visite des malades et auront, chemin faisant, la plus grande attention sur ceux qui seraient menacés de la maladie, pour les en garantir, s’il était possible, en leur prescrivant la diète, les boissons les émétique, les purgatifs et les stomachiques convenables.

\medskip

\textbf{4/} Vu l’état de dépérissement et de faiblesse dans lequel se trouvent la plupart des malades ils soutiendront leurs forces par de bons bouillons, de légers cordiaux et des tisanes légèrement incisives….Ils seront très réservés sur les saignées, les émétique et les purgatifs… en espérant que les approches d’une saison plus tempérée contribuera à terminer bientôt cette fâcheuse maladie.

\vspace{0.35em}

\begin{tcolorbox}[
  colback=white,
  colframe=black,
  boxrule=0.7pt,
  arc=0pt,
  left=6pt,right=6pt,top=5pt,bottom=5pt
]
\textbf{On peut penser, sans grand risque d’erreur, que la misère extrême du pays et particulièrement des vallées comme du village de Lanuéjols, autrefois protestant, était une conséquence des persécutions et des dragonnades qui, après le brûlement des Cévennes de 1703, continuèrent tout au long du 18\ieme{} siècle.}
\end{tcolorbox}

}


